% BHU PYQ -- Professional Exam Template
\documentclass[11pt,a4paper]{article}

\usepackage[a4paper,top=2.5cm,bottom=2.5cm,left=2.8cm,right=2.8cm,headheight=14pt]{geometry}
\usepackage{amsmath,amssymb}
\usepackage[T1]{fontenc}
\usepackage[utf8]{inputenc}
\usepackage{lmodern}
\usepackage{microtype}
\usepackage{fancyhdr}
\usepackage{enumitem}
\usepackage{hyperref}
\usepackage{lastpage}
\usepackage{parskip}

\hypersetup{colorlinks=true,urlcolor=black,linkcolor=black,
  pdftitle={B.Sc. (Hons.) Botany Sem-VI BOB-604 -- BHU PYQ}}

\pagestyle{fancy}
\fancyhf{}
\renewcommand{\headrulewidth}{0.4pt}
\renewcommand{\footrulewidth}{0.4pt}
\fancyhead[L]{\small   \textit{Banaras Hindu University}}
\fancyhead[R]{\small   \textit{Botany Paper: BOB-604}}
\fancyfoot[C]{\small Page  \thepage\ of \pageref{LastPage}}


\newlist{parts}{enumerate}{2}
\setlist[parts,1]{label=(\alph*),leftmargin=*,itemsep=4pt,topsep=3pt}
\setlist[parts,2]{label=(\roman*),leftmargin=*,itemsep=2pt,topsep=2pt}

\newcommand{\pts}[1]{\hfill   \textit{[#1]}}


\begin{document}


\begin{center}
  {\large\bfseries BANARAS HINDU UNIVERSITY}\\[2pt]
  {
\normalsize B.Sc. (Hons.) Semester VI Examination, 2022-23}\\[6pt]
  \rule{0.8\linewidth}{0.5pt}\\[6pt]
  {\large\bfseries Botany}\\[4pt]
  {\large\bfseries Paper: BOB-604 --- Plant and Microbial Techniques}\\[4pt]
  {
\normalsize Time: Three Hours \hfill Maximum Marks: 70}
\end{center}

\rule{\linewidth}{0.4pt}

\smallskip



\noindent   \textit{Instructions: Answer five questions including Question 1 which is compulsory. The figures in the right-hand margin indicate marks.}

\bigskip

\medskip


\noindent   \textbf{Question 1.} Answer the following in not more than 50 words each:\pts{$10  \times 1 = 10$}

\begin{parts}
  \item What is the role of TCA in separation of protein?
  \item What is EMBL?
  \item What is DEAE and ion-exchange chromatography?
  \item What is DEPC?
  \item What is the function of Ribonuclease H (RNase H)?
  \item What are Cy3 and Cy5?
  \item What is the difference between magnification and resolving power of a lens?
  \item What is Stokes shift?
  \item Discuss the surface sterilant in tissue culture.
  \item What is protoplast fusion?
\end{parts}

\medskip


\noindent   \textbf{Question 2.} Answer any two of the following:\pts{$2  \times 7.5 = 15$}

\begin{parts}
  \item Describe the principle of operation and application of gas chromatography.
  \item Illustrate the process of PCR with a suitable diagram.
  \item What is gel filtration chromatography? Explain its merits and demerits.
\end{parts}

\medskip


\noindent   \textbf{Question 3.} What is a Transmission Electron Microscope (TEM)? Describe the principle and components of TEM with a suitable diagram.\pts{15}

\medskip


\noindent   \textbf{Question 4.} Write short notes on any two of the following:\pts{$2  \times 7.5 = 15$}

\begin{parts}
  \item Components and their specific functions of a plant tissue culture medium
  \item Fluorescence microscopy
  \item DNA microarray
\end{parts}

\medskip


\noindent   \textbf{Question 5.} What is micropropagation? Explain the different stages involved in the process with suitable diagrams. Describe its advantages and limitations.\pts{15}

\medskip


\noindent   \textbf{Question 6.} What is a cDNA library? How is a cDNA library prepared and what are its applications?\pts{15}

\medskip


\noindent   \textbf{Question 7.} Describe any two of the following:\pts{$2  \times 7.5 = 15$}

\begin{parts}
  \item Principle and applications of ion-exchange chromatography
  \item Difference between selective media and indicator media with examples
  \item Principle of operation and general scheme of a mass spectrometer
\end{parts}

\medskip


\noindent   \textbf{Question 8.} What is automated DNA sequencing? Describe in detail (a) ILLUMINA and (b) Roche 454, along with their merits and demerits.\pts{15}

\vfill
\rule{\linewidth}{0.4pt}

\begin{center}\small
\href{https://bhu-exam.vercel.app/}{Click Here}\\[4pt]\href{https://me-aryan.vercel.app/}{Click Here}\\[4pt]\href{https://quashan-me.vercel.app/}{Click Here}
\end{center}

\end{document}
